\documentclass[12pt]{article}
\usepackage[margin=1in]{geometry}

% Start of preamble
%==========================================================================================%
% Required to support mathematical unicode
\usepackage[warnunknown, fasterrors, mathletters]{ucs}
\usepackage[utf8x]{inputenc}

\usepackage[dvipsnames,table,xcdraw]{xcolor}
\usepackage{hyperref} 
\hypersetup{
colorlinks=true,
linkcolor=blue,
filecolor=magenta,
urlcolor=cyan,
pdfpagemode=FullScreen
}

% Standard mathematical typesetting packages
\usepackage{amsmath,amssymb,amscd,amsthm,amsxtra, pxfonts}
\usepackage{mathtools,mathrsfs,dsfont,xparse}

% Symbol and utility packages
\usepackage{cancel, textcomp}
\usepackage[mathscr]{euscript}
\usepackage[nointegrals]{wasysym}
\usepackage{apacite}

% Extras
\usepackage{physics}  
\usepackage{tikz-cd} 
\usepackage{microtype}
\usepackage{enumitem}
\usepackage{titling}
\usepackage{graphicx}

% Fancy theorems due to @intuitively on discord
\usepackage{mdframed}
\newmdtheoremenv[
backgroundcolor=NavyBlue!30,
linewidth=2pt,
linecolor=NavyBlue,
topline=false,
bottomline=false,
rightline=false,
innertopmargin=10pt,
innerbottommargin=10pt,
innerrightmargin=10pt,
innerleftmargin=10pt,
skipabove=\baselineskip,
skipbelow=\baselineskip
]{mytheorem}{Theorem}

\newenvironment{theorem}{\begin{mytheorem}}{\end{mytheorem}}

\newtheorem{corollary}{Corollary}
\newtheorem{lemma}{Lemma}

\newtheoremstyle{definitionstyle}
{\topsep}%
{\topsep}%
{}%
{}%
{\bfseries}%
{.}%
{.5em}%
{}%
\theoremstyle{definitionstyle}
\newmdtheoremenv[
backgroundcolor=Violet!30,
linewidth=2pt,
linecolor=Violet,
topline=false,
bottomline=false,
rightline=false,
innertopmargin=10pt,
innerbottommargin=10pt,
innerrightmargin=10pt,
innerleftmargin=10pt,
skipabove=\baselineskip,
skipbelow=\baselineskip,
]{mydef}{Definition}
\newenvironment{definition}{\begin{mydef}}{\end{mydef}}

\newtheorem*{remark}{Remark}

\newtheorem*{example}{Example}

% Common shortcuts
\def\mbb#1{\mathbb{#1}}
\def\mfk#1{\mathfrak{#1}}

\def\bN{\mbb{N}}
\def \C{\mbb{C}}
\def \R{\mbb{R}}
\def\bQ{\mbb{Q}}
\def\bZ{\mbb{Z}}
\def \cph{\varphi}
\renewcommand{\th}{\theta}
\def \ve{\varepsilon}
\newcommand{\mg}[1]{\| #1 \|}

% Often helpful macros
\newcommand{\floor}[1]{\left\lfloor#1\right\rfloor}
\newcommand{\ceil}[1]{\left\lceil#1\right\rceil}
\renewcommand{\qed}{\hfill\qedsymbol}
\renewcommand{\P}{\mathbb P\qty}
\newcommand{\E}{\mathbb{E}\qty}
\newcommand{\Cov}{\mathrm{Cov}\qty}
\newcommand{\Var}{\mathrm{Var}\qty}

% Sets
\usepackage{braket}

\graphicspath{{/}}
\usepackage{float}

\newcommand{\SET}[1]{\Set{\mskip-\medmuskip #1 \mskip-\medmuskip}}

% End of preamble
%==========================================================================================%

% Start of commands specific to this file
%==========================================================================================%

%==========================================================================================%
% End of commands specific to this file

\title{CSE Template}
\date{\today}
\author{Rohan Mukherjee}

\begin{document}
    \maketitle
    \subsection*{A1.}  
    \begin{enumerate}[label=\alph*.]
        \item In the L1 norm, small numbers are punished more than in the L2 norm. For example, in the L2 norm a coordinate of 0.1 translates to an increase in only 0.01, which is almost negligible. But in the L1 norm it still has a lot of weight.
        \item This is true, since we might step past the minimum with too large a step size, and then we will try to step back towards it, but just end up oscillating far from the minimum.
        \item SGD runs much, much faster than GD, bringing the running time from $O(n \cdot d)$ per iteration to only $O(d)$. GD on the other hand is much more stable, not relying on randomness and we don't have to worry about conrtrolling the step size.
        \item Logistic regression does not have a closed form solution, while linear regression does.
    \end{enumerate}

    \subsection*{A2.}
    \begin{enumerate}
        \item Notice that,
        \begin{align*}
            (|a+b|)^2 = (a+b)^2 = a^2+b^2+2ab \leq a^2+b^2+2|a||b| = (|a|+|b|)^2
        \end{align*}
        Since $f(x) = \sqrt(x)$ is increasing, and both the LHS and RHS are positive, taking square roots on both sides shows that $|a+b| \leq |a| + |b|$.
        \item Consider the vectors $x = (1/4, 0)$ and $y = (0, 1/4)$. Then,
        \begin{align*}
            g(x+y) = (\frac{1}{2} + \frac{1}{2})^2 = 1
        \end{align*}
        while,
        \begin{align*}
            g(x) + g(y) = \frac{1}{4} + \frac{1}{4} = \frac{1}{2}
        \end{align*}
        Which shows that the triangle inequality does not hold.
    \end{enumerate}

    \subsection*{A3.}
    Note that the set of points $\SET{\lambda a + (1-\lambda)b \mid 0 \leq \lambda \leq 1} = \SET{a + \lambda(b-a) \mid 0 \leq \lambda \leq 1}$ is precisely the lien conneting $a$ with $b$. Since in part I, the line connecting the point $b$ with $c$ does not live fully within the set, we see that I is not convex. Similarly, applying the same logic to points a and d in the second picture shows that the II is not convex either.

    \subsection*{A4.}
    The first function I is convex because it is a quadratic function pointing upwards, with always-positive second derivative. The second function is not convex because The secant line between $b$ and $c$ lies below the curve, not above it (Recall that the line $\lambda f(x) + (1-\lambda)f(y)$ for $0 \leq \lambda \leq 1$ is the secant line of $f$ between $x$ and $y$).

    \subsection*{A7.}
    \begin{enumerate}[label=\alph*.]
        \item Notice that
        \begin{align*}
            \dv{x} \log(1+e^{-x}) = \frac{-e^{-x}}{1+e^{-x}} = \frac{1}{1+e^{-x}}-1
        \end{align*}
        We recall the chain rule: if $f: \R \to \R$, $g: \R^n \to \R$ are differentiable, then $\nabla_x (f \circ g)(x) = f'(g(x)) \nabla_x g(x)$. Taking $f(x) = \log(1+e^{-x})$ and $g(x) = y_i(b+x_i^Tw)$, we get that
        \begin{align*}
            \nabla_w f(g(x)) = \qty(\frac{1}{1+e^{-y_i(b+x_i^Tw)}}-1) y_ix_i = (\mu_i(w,b) - 1)y_ix_i
        \end{align*}
        Because $\nabla_w y_i(b+x_i^Tw) = \nabla_w y_ix_i^Tw = y_ix_i$. Similarly, 
        \begin{align*}
            \pdv{b} f(g(x)) = \qty(\frac{1}{1+e^{-y_i(b+x_i^Tw)}}-1) y_i = (\mu_i(w,b)-1)y_i
        \end{align*}
        Since $\grad_w \lambda w^Tw = \lambda 2w$, we have that
        \begin{align*}
            \grad_w J(w,b) &= \frac1n\sum_{i=1}^n \qty(\mu_i(w,b)-1)y_ix_i + 2\lambda w \\
            \pdv{b} J(w,b) &= \frac1n\sum_{i=1}^n \qty(\mu_i(w,b)-1)y_i
        \end{align*}
    \end{enumerate}
\end{document}